
\section{Amplificador no inversor}

La configuración no inversora es similar a la configuración inversora, sin embargo en este circuito la señal de salida no se invierte (en fase) respecto de la entrada.
\begin{figure}[!ht]
  \centering
  \begin{circuitikz}
    \node[op amp](OA) at (0,0){};
    \draw (OA.-) node[left]{\(+\)} to[short,-*] ++(0,1) coordinate(tmp) to[R,l_={\(R_1\)},i<_={\(i_1\)}] ++(-3,0) node[ground,rotate=-90]{};
    \draw (tmp) to[R,l={\(R_2\)},i={\(i_2\)}] (tmp -| OA.out) to[short,-*] (OA.out) to[short,-o] ++(1,0) node[right]{\(v_\text{out}\)};
    \draw (OA.+) to[short] ++(0,-.5) to[short,-o] ++(-1,0) node[left]{\(v_\text{in}\)};
    \node[left] at (OA.+){\(-\)};
    \node[left] at(-1,0){\(v_i\)};
  \end{circuitikz}
  \caption{Diagrama topológico de la configuración no inversora.}
\end{figure}

Del mismo modo que en la configuración inversora, como la corriente entrante al amplificador es cero, \(i_1=i_2\):
\begin{align*}
  i_1&=i_2\\ 
  \frac{0-v_i}{R_1} &= \frac{v_i-v_\text{out}}{R_2}
\end{align*}
Pero como \(v_i=v_\text{in}\),
\begin{align*}
  -\frac{v_\text{in}}{R_1} &= \frac{v_\text{in}}{R_2} - \frac{v_\text{out}}{R_2}\\ 
  v_\text{in}\left(\frac{1}{R_1}+\frac{1}{R_2}\right) &= v_\text{out}\frac{1}{R_2}
\end{align*}
Reordenando, resulta
\begin{align*}
  \frac{v_\text{out}}{v_\text{in}} &= \frac{R_2}{R_1} + \frac{R_2}{R_2}\\
  \frac{v_\text{out}}{v_\text{in}} &= \frac{R_2}{R_1}+1
\end{align*}
Entonces, la amplificación es:
\begin{equation}\label{eq_ampl_no_inv}
  \boxed{A=1+\frac{R_2}{R_1}}
\end{equation}

\subsection{Seguidor de tensión}

Si \(R_1\to\infty\), es decir, el terminal inversor aisla completamente respecto de tierra, y \(R_2\) se ignora (ya que si se observa la expresión \eqref{eq_ampl_no_inv} si \(R_1\to\infty\) luego \(A=1\) sin importar el valor de \(R_2\)), se obtiene un seguidor de tensión. El diagrama topológico de este circuito es como se muestra en la figura \ref{fig_seguidor}.
\begin{figure}[!ht]
  \centering
  \begin{circuitikz}
    \node[op amp](A){};
    \draw (A.-) to[short] ++(0,1) coordinate(tmp) to[short] (tmp -| A.out) to[short,-*] (A.out) to[short,-o] ++(1,0) node[right]{\(v_\text{out}\)};
    \draw (A.+) to[short] ++(0,-.5) to[short,-o] ++(-.5,0) node[left]{\(v_\text{in}\)};
  \end{circuitikz}
  \caption{Configuración de seguidor de tensión.}
  \label{fig_seguidor}
\end{figure}

Este tipo de circuito es muy utilizado como etapa separadora entre circuitos con relativa alta impedancia de salida y de entrada.

